% =====================================================================
%  §2 — Background and Related Work
%
%  Seven families of standards: FHIR Evidence, ECO, SEPIO, GA4GH GKS,
%  ACMG-AMP, IAO/OBI, PROV-O. Each in a compact paragraph. Closes
%  with a paragraph naming the gap this paper addresses.
%
%  Revised 2026-04-21: standards-mapping table dropped for space
%  (paper was at 18pp vs 12pp cap); per-family paragraphs tightened
%  and acronyms expanded consistently on first mention.
% =====================================================================

\section{Background and Related Work}
\label{sec:supp-background}

Scientific and clinical evidence has been the subject of sustained standards
work in clinical informatics, biomedical ontology, and genomic knowledge
representation. Our schema lies at the intersection of these communities. It
is oriented toward downstream decision-making, grounded in formal constraints
and explicit provenance, and treats target resolution, including variant,
gene, region, and whole-genome aggregate, as a first-class concern while
relying on GA4GH standards for sequence identifiers.

This note positions the schema relative to seven families of standards and
identifies the coverage gap that motivated the work. A first term-level
alignment is provided in \Cref{sec:supp-crosswalk}. Completing the coverage
and formalizing the mappings in OWL remain future work.

\paragraph{FHIR Evidence Resources.}
The HL7 FHIR Evidence, EvidenceVariable, and EvidenceReport
resources~\cite{fhir_evidence} provide structured representations of
clinical-trial-style evidence, with explicit fields for population,
intervention, comparator, outcome, and statistical measures. They were
designed with systematic-review and guideline-development workflows in mind,
where evidence concerns an intervention tested in a population. This framing
does not fit basic-science and pre-clinical genetic evidence cleanly. A mouse
knockout study, for example, does not have an ``intervention'' in the
clinical-trial sense, and a coding-variant case--control study has a target,
the variant, that is not the same kind of object as a drug or procedure. Our
schema shares FHIR Evidence's commitment to structured, dimensional
description but replaces its clinical-trial frame with a vocabulary suited to
genetic claims.

\paragraph{Evidence Codes: ECO and SEPIO.}
The Evidence and Conclusion Ontology (ECO)~\cite{eco2019} provides a large
hierarchical vocabulary of evidence types, including inference methods, assay
types, and source-document categories. It is organized within the Open
Biological and Biomedical Ontologies (OBO) framework and is widely used in
Gene Ontology annotation. The Scientific Evidence and Provenance Information
Ontology (SEPIO)~\cite{sepio} complements ECO with machine-readable
representations of statements, qualifiers, and provenance relationships. Our
Method and Measurement Target dimensions overlap substantially with areas
already covered by ECO and SEPIO. Our long-term goal is to ground those
dimensions in ECO and SEPIO identifiers rather than maintain a separate
vocabulary.

\paragraph{GA4GH Genomic Knowledge Standards.}
The most relevant comparison within the GA4GH Genomic Knowledge Standards
effort is the Variant Annotation (VA) Specification~\cite{ga4gh_va}. VA is a
broad information model for structured statements about variants, organized
around a \texttt{Statement} class whose instances carry a proposition,
direction of support, strength, classification, evidence lines, and
contribution metadata. VA is information-centric: its top-level abstraction,
\texttt{InformationEntity}, represents what is known about a variant rather
than the process by which primary-literature claims become evidence. Our
schema is narrower and complementary. It focuses on evidence drawn from the
primary literature, with the \emph{Relevance and Specificity} and
\emph{Confidence} challenges identified in the main-text introduction as
concrete targets. The Variation Representation Specification
(VRS)~\cite{ga4gh_vrs}, a separate product of the same effort, addresses
canonical variant identity across coordinate systems and is orthogonal to the
present work.

\paragraph{ACMG/AMP variant interpretation.}
The American College of Medical Genetics and Genomics and Association for
Molecular Pathology (ACMG/AMP) guidelines~\cite{acmg2015} specify a
rule-based framework for variant classification from weighted evidence codes
and criteria. Our schema is intended to supply inputs to such frameworks,
rather than replace them: its per-publication evidence items are the kind of
input an ACMG/AMP-compatible classifier would consume. A term-level alignment
of the Method and Credibility dimensions with ACMG/AMP evidence codes remains
future work.

\paragraph{Upper ontologies: IAO and OBI.}
The Information Artefact Ontology (IAO)~\cite{iao2015} and the Ontology for
Biomedical Investigations (OBI)~\cite{obi2016} provide upper-level
scaffolding for many OBO-compliant biomedical ontologies. IAO supplies
classes such as \emph{information-content-entity} and
\emph{textual-entity}; OBI supplies assay and protocol hierarchies relevant
to the Method dimension. Formal alignment of the model's top-level classes
with IAO and OBI, together with an OWL/LinkML rendering of the schema,
remains future work.

\paragraph{Provenance: PROV-O.}
The W3C PROV Ontology (PROV-O)~\cite{prov_o} provides standard provenance
relationships. Aligning the model with PROV-O remains future work.

\paragraph{Related annotation efforts.}
The Clinical Genome Resource (ClinGen) evidence-curation
framework~\cite{clingen2018} provides a structured workflow for gene--disease
validity and variant-pathogenicity assessments with explicit evidence
rubrics. Our schema shares its commitment to source-anchored claims and
explicit categorization, but operates at a lower level by representing the
individual claims that such frameworks aggregate. The Monarch Disease Ontology
(MONDO)~\cite{mondo} and the Human Phenotype Ontology (HPO)~\cite{hpo}
provide established vocabularies for disease and phenotype terms. We expect
curator-authored HPO terms to populate the schema's phenotype-mapping fields,
as illustrated by the HPO-term suggestions in the Davis, Nelson, and Gupta
case studies (Section~5 of the main paper).

\subsection{Requirements and coverage matrix}
\label{sec:supp-coverage-matrix}

The per-family paragraphs above describe each standard on its own terms.
\Cref{tab:coverage-matrix} takes the complementary perspective: it begins
with ten evidence patterns that a genetic publication may report and asks
whether each pattern can be represented by each standard and by GEM. Eight
patterns are exercised by the six-paper corpus in the main paper
(\Cref{tab:corpus}). The final two, therapeutic-response evidence and
cross-publication synthesis, were exposed by the corpus as patterns outside
the current model's scope.

The matrix is a coverage inventory, not a ranking. The standards were
developed for different purposes, and a pattern marked as not represented is
a limitation only if that pattern lies within the standard's intended scope.

The columns represent the five standards families in the comparison that
carry evidence content. IAO, OBI, and PROV-O are omitted because they are
orthogonal to all ten patterns. They provide upper-level classes, assay
hierarchies, and provenance relations that any of the represented models
could reuse, but do not themselves represent the patterns.

Cell values are interpreted relative to each standard's role. For an
information model, including FHIR Evidence, GA4GH VA, and SEPIO,
\cfull{} indicates that a class or published profile represents the
pattern. \chalf{} indicates that the pattern is representable only through
the standard's extension mechanism, such as a custom VA profile, a FHIR
extension, or an external predicate vocabulary, or that the standard
represents only one aspect of the pattern. \cnone{} indicates that the
pattern lies outside the model's frame. For ECO, \cfull{} indicates that a
term for the pattern exists, \chalf{} that only a generic ancestor applies,
and \cnone{} that no suitable term was identified.

For the ClinGen gene--disease validity framework~\cite{strande2017} and
the ACMG/AMP variant-interpretation guidelines~\cite{acmg2015},
\cfull{} indicates that the pattern is a scored evidence category.
\chalf{} indicates partial or indirect treatment, and \cnone{} indicates
that neither framework provides a corresponding scored category. For GEM,
\cfull{} indicates that the pattern is represented natively and was
annotated in the corpus without a forced fit. \chalf{} indicates that the
pattern was annotated only through a forced fit recorded as a candidate
extension, identified by a CE code in
\fpath{schema/EXTENSIONS.md}. \cnone{} indicates that the pattern is not
represented. Every identifier cited in the row notes was verified against
the relevant source vocabulary or specification at the time of writing.

\begin{table}[H]
  \centering
  \footnotesize
  \caption[Requirements and coverage matrix]{Ten evidence patterns across
  five families of standards and GEM. \cfull{} represented natively;
  \chalf{} represented only through the standard's extension mechanism, or
  only in part; \cnone{} not represented; \cna{} outside the standard's
  concern. For GEM, \chalf{} denotes a forced fit recorded as a candidate
  extension. The \emph{remaining gap} column is relative to GEM and names
  the candidate extension where one exists. Row notes follow the table.
  FHIR = HL7 FHIR Evidence R5; VA = GA4GH Variant Annotation 1.0 with
  VRS; ClinGen/ACMG = the ClinGen gene--disease validity framework and the
  ACMG/AMP variant-interpretation guidelines.}
  \label{tab:coverage-matrix}
  \setlength{\tabcolsep}{4pt}
  \begin{tabularx}{\linewidth}{@{}>{\raggedright\arraybackslash}p{0.30\linewidth} c c c c c c >{\raggedright\arraybackslash}X@{}}
    \toprule
    Evidence pattern & FHIR & VA & SEPIO & ECO
      & \begin{tabular}[b]{@{}c@{}}ClinGen/\\ACMG\end{tabular} & GEM
      & Remaining gap (GEM) \\
    \midrule
    1. Mendelian pedigree: variant co-segregates with disease in a family
      & \chalf & \cfull & \chalf & \cnone & \cfull & \cfull
      & segregation strength (meioses, LOD) unstructured \\
    2. Case--control burden: gene-level aggregate of rare variants
      & \chalf & \chalf & \chalf & \chalf & \cfull & \cfull
      & cohort composition (CE-IN5) \\
    3. Common-variant association: GWAS discovery, replication, fine mapping
      & \chalf & \chalf & \chalf & \chalf & \cnone & \cfull
      & direction of effect (CE-DU1); activation scope (CE-DU4); no discovery--replication link \\
    4. Functional assay: variant or gene effect measured \emph{in vitro}
      & \chalf & \cfull & \chalf & \cfull & \cfull & \cfull
      & sub-gene resolution (CE-J1); assay type is free text \\
    5. Animal model: knockout or conditional-knockout phenotype
      & \chalf & \chalf & \chalf & \cfull & \cfull & \cfull
      & HPO-only phenotype mapping; rescue not distinguished \\
    6. Gene--gene relation: X inhibits, binds, or regulates Y
      & \cnone & \chalf & \chalf & \cfull & \chalf & \chalf
      & relation vocabulary (CE-J3); complex as target (CE-J2) \\
    7. Variant identity across coordinate systems (protein, cDNA, genomic)
      & \chalf & \cfull & \cna & \cna & \cfull & \chalf
      & cross-level description (CE-N1, CE-C1); VRS/HGVS binding \\
    8. Polygenic score: construction and validation of a genome-wide aggregate
      & \chalf & \chalf & \chalf & \chalf & \chalf & \chalf
      & score machinery (CE-IN1 to CE-IN6) \\
    9. Therapeutic response: intervention tested in patients
      & \cfull & \cfull & \chalf & \cfull & \chalf & \cnone
      & no intervention construct (CE-DU3) \\
    10. Cross-publication synthesis: evidence lines combined into a classification
      & \cfull & \cfull & \cfull & \chalf & \cfull & \cnone
      & out of scope by design (CE-DU3) \\
    \bottomrule
  \end{tabularx}
\end{table}

\paragraph{Row notes.}
Paper keys follow \Cref{tab:corpus}; item identifiers (GE-$n$) follow the
released annotations.

\begin{enumerate}
  \item \emph{Pedigree segregation.} FHIR can represent the family as a
    \texttt{Group} population and the segregation statistic as an ad hoc
    \texttt{statistic}, but it has no dedicated segregation construct. VA
    covers this pattern through the ACMG 2015 Variant Pathogenicity Evidence
    Line profile, in which co-segregation (PP1) is a criterion that an
    evidence line may carry. SEPIO provides an assertion and evidence-line
    frame, but the segregation content must come from a domain model, as in
    the SEPIO-based ClinGen interpretation model. No ECO term for
    co-segregation, pedigree, or linkage evidence was identified. The only
    retrieved term for ``segregation'' was the generic ECO:0005613,
    \emph{inference by association of genotype from phenotype}. ClinGen
    scores segregation as case-level evidence~\cite{strande2017}, and
    ACMG/AMP codes it as PP1~\cite{acmg2015}. GEM represents the pattern
    natively through Mode of Inheritance and Mendelian Segregation, as
    exercised by Gupta GE-1 (X-linked recessive) and Davis GE-5
    (autosomal recessive). Because Mendelian Segregation is Boolean,
    segregation strength is recorded only in the credibility facets.

  \item \emph{Case--control burden.} FHIR represents the study design and
    statistic natively through \texttt{studyDesign} and \texttt{statistic},
    but the gene-level aggregate exposure must be coded ad hoc. VA 1.0 has
    no gene-level proposition among its base profiles. Its nearest study
    result, Cohort Allele Frequency, is variant-level, so a custom profile
    is required. ECO offers only the generic ECO:0005613. ClinGen scores
    case--control evidence, including aggregate variant
    analysis~\cite{strande2017}. GEM represents the pattern natively in
    Davis GE-1, with Population Genetics and Human Genetics, Statistical
    Genetics, and a Gene target. Case and control counts remain captured as
    a credibility facet, motivating CE-IN5.

  \item \emph{Common-variant association.} FHIR and VA have the same
    limitations described in row 2. Neither distinguishes a genome-wide
    discovery stage from a replication stage, and VA has no association
    proposition. ECO has no genome-wide-association or replication term
    beyond ECO:0005613. The ClinGen validity framework scores Mendelian
    gene--disease evidence and does not score complex-trait association.
    GEM represents the pattern natively through the Method hierarchy:
    GWAS, Association Study, and Fine Mapping under Statistical Genetics.
    This is exercised by Duerr GE-1, GE-2, and GE-4. The same paper exposed
    the missing direction-of-effect dimension for a protective allele
    (CE-DU1) and, through the family-based transmission test in GE-3, an
    activation gap for Human Genetics with a Variant target (CE-DU4).

  \item \emph{Functional assay.} FHIR has no population--intervention frame
    for a cell-based assay; the assay system would have to be cast as a
    population. VA covers the variant-level case through the Experimental
    Variant Functional Impact proposition and study-result profiles, but
    gene-level assays require a custom profile. ECO:0000181,
    \emph{in vitro assay evidence}, and its descendants cover the evidence
    type. ClinGen scores experimental evidence, including biochemical
    function and functional alteration~\cite{strande2017}; ACMG/AMP codes
    it as PS3 or BS3. GEM represents this pattern natively, as illustrated
    by Nelson GE-1 to GE-3 and the \emph{in vitro} Jossin items. Jossin
    exposed the missing protein-subdomain resolution (CE-J1). The assay
    itself is described only at the Method level, \emph{In Vitro}, pending
    planned ECO/OBI grounding (\cref{sec:supp-crosswalk}).

  \item \emph{Animal model.} FHIR can cast the genotype as an exposure and
    the animal cohort as a \texttt{Group}, but it has no model-organism
    frame. VA has no animal-model proposition; the study can enter only as
    an evidence item behind a custom profile. ECO provides detailed
    evidence-type coverage through ECO:0000179,
    \emph{animal model system study evidence}; ECO:0001091,
    \emph{knockout phenotypic evidence}; ECO:0001030,
    \emph{conditional knockout evidence}; and ECO:0000013,
    \emph{transgenic rescue experiment evidence}. ClinGen scores
    non-human model organisms and rescue as experimental
    evidence~\cite{strande2017}. GEM represents the pattern natively with
    the conditional Organism and Knockout Type dimensions. Examples include
    Jossin GE-1, a conditional mouse knockout with organismal-scale
    evidence, and Davis GE-2 in zebrafish and GE-4 in rat. The
    phenotype-mapping fields currently expect HPO terms, so they do not bind
    model-organism phenotype vocabularies, and rescue is not distinguished
    from other \emph{in vivo} evidence.

  \item \emph{Gene--gene relation.} FHIR Evidence has no frame for a
    relation between two molecular entities. VA includes a \texttt{Gene}
    domain entity but no interaction proposition. A SEPIO assertion can
    carry the relation through a predicate from the Relation Ontology.
    ECO:0000021, \emph{physical interaction evidence}, covers the evidence
    type. ClinGen counts protein interaction only as supporting
    experimental evidence~\cite{strande2017}, not as a claim type. GEM has
    a Gene Relation dimension, of whose three current values the corpus
    exercised one: X inhibits Y in Jossin GE-5. Binding and trafficking
    relations required by other Jossin items were forced fits (CE-J3), as
    was a protein complex as a target (CE-J2).

  \item \emph{Variant identity across coordinate systems.} This pattern is
    outside the FHIR Evidence module. FHIR handles it through genomics
    reporting resources that an \texttt{Evidence} resource can reference.
    VA delegates to VRS~\cite{ga4gh_vrs}, which was designed for this
    purpose. SEPIO and ECO do not address variant identity. ClinGen resolves
    identity through HGVS and the ClinGen Allele Registry. GEM relies on
    GA4GH identifiers by design, but its Resolution values currently
    presuppose genomic coordinates. Protein-level descriptions in Nelson
    (CE-N1) and the transcript-level pattern identified during the UMLS
    crosswalk (CE-C1) were forced fits. Binding target identifiers to VRS or
    HGVS is planned rather than implemented.

  \item \emph{Polygenic score.} FHIR can represent a risk-prediction study
    with the score as a coded variable, but it has no dedicated score
    construct. VA 1.0 and VRS provide no entity for a genome-wide aggregate,
    so a custom profile would be required. ECO offers only the generic
    computational term ECO:0007672 and association term ECO:0005613.
    ClinGen's relevant contribution is the PRS-RS reporting
    checklist~\cite{wand2021}, not a curation rubric. In GEM, all four
    Inouye items are forced fits: the target is recorded as Variant, the
    Resolution value is a placeholder, and Variant Ascertainment is set to
    the documented sentinel. These conditions motivated CE-IN1 through
    CE-IN6, which await a second polygenic-score paper under the
    two-papers rule.

  \item \emph{Therapeutic response.} This pattern fits FHIR Evidence's
    native population--intervention--comparator--outcome frame. VA covers
    it through the Variant Therapeutic Response proposition. ECO:0000180,
    \emph{clinical study evidence}, covers the evidence type. ClinGen
    addresses therapeutic response through actionability curation rather
    than the gene--disease validity framework. GEM has no intervention
    construct. The cited anti-p40 antibody trial in Duerr, GE-new-1, was
    recorded with Knowledge Domain left unresolved by curator decision rather
    than forced into a genetic domain. It is the second independent instance
    of CE-DU3.

  \item \emph{Cross-publication synthesis.} Combining evidence into a
    conclusion is the purpose of FHIR \texttt{EvidenceReport}; of VA
    \texttt{Statement}s with evidence lines, including ACMG classification
    profiles; of SEPIO evidence lines; and of the ClinGen and ACMG/AMP
    rubrics. ECO labels the result only as an evidence type through
    ECO:0000212, \emph{combinatorial evidence}. GEM does not synthesize
    evidence by design. It represents the per-publication evidence items
    consumed by such frameworks, while the reasoning layer that combines
    them is developed separately~\cite{our_preprint}.
\end{enumerate}

Read by column, the matrix clarifies GEM's contribution. Rows 1 through 5
cover common patterns in basic-science genetic publications. Among the
columns represented here, GEM is the only one that represents all five
natively within a single dimensional frame. Other standards either require
an extension mechanism because they were designed for a different primary
purpose, or provide an evidence-type vocabulary that names the evidence
without structuring the claim. Rows 6 through 8 identify patterns that the
pilot exposed as forced fits and therefore as documented areas for extension.
Rows 9 and 10 are outside GEM's intended scope, and the matrix indicates
which standards can represent them.

\paragraph{The coverage gap.}
None of these standards was designed to treat the heterogeneous evidence
patterns of basic-science and pre-clinical genetic publications as a
first-class concern. A Llgl1 conditional-knockout study, for example, may
report molecular, cellular, and clinical phenotypes in the same paper. A
polygenic-score study may aggregate millions of variants into one score. Each
existing standard captures part of what such papers report, but none captures
the whole. The schema introduced in the main paper aims to represent this
heterogeneity in a form that is compact enough for manual annotation,
machine-checkable enough for automated validation, and flexible enough to
accommodate the evidence patterns produced by the literature.
